\documentclass[12pt]{article}

\usepackage{amsmath,amssymb,amsthm,mathrsfs}
\usepackage{physics}
\usepackage{hyperref}
\usepackage{geometry}
\geometry{margin=1in}

\title{\textbf{Spectral Geometry and Limiting Structures of the Morse Hamiltonian}}
\author{ }
\date{}

\newtheorem{theorem}{Theorem}[section]
\newtheorem{proposition}[theorem]{Proposition}
\newtheorem{lemma}[theorem]{Lemma}
\newtheorem{definition}[theorem]{Definition}
\newtheorem{remark}[theorem]{Remark}

\begin{document}

\maketitle

\begin{abstract}
We present a complete spectral-theoretic analysis of the Morse Hamiltonian on the real line, including bound states, scattering states, self-adjointness, deficiency indices, resonances, and limiting regimes.
Special attention is given to the hard-wall limit $a\to\infty$ and its comparison with large-parameter limits of the P\"oschl--Teller and Scarf families.
The analysis clarifies how boundary formation, algebraic contraction, and hypergeometric confluence govern spectral transitions.
\end{abstract}

\tableofcontents

%%%%%%%%%%%%%%%%%%%%%%%%%%%%%%%%%%%%%%%%%%%%%%%%%%%%%%%%%%%%
\section{The Morse Hamiltonian}

We consider
\[
H = -\frac{\hbar^2}{2m}\frac{d^2}{dx^2}
+ D_e(1-e^{-a(x-x_0)})^2,
\quad x\in\mathbb R,
\]
with parameters $D_e,a>0$.

Define
\[
\lambda=\frac{\sqrt{2mD_e}}{\hbar a}.
\]

%%%%%%%%%%%%%%%%%%%%%%%%%%%%%%%%%%%%%%%%%%%%%%%%%%%%%%%%%%%%
\section{Self-Adjointness}

\begin{theorem}
The Morse Hamiltonian is essentially self-adjoint on $C_c^\infty(\mathbb R)$.
\end{theorem}

\begin{proof}
As $x\to -\infty$, $V(x)\to +\infty$.
As $x\to +\infty$, $V(x)\to D_e$.
The operator is limit-point at both endpoints by Weyl’s alternative.
Thus deficiency indices vanish.
\end{proof}

%%%%%%%%%%%%%%%%%%%%%%%%%%%%%%%%%%%%%%%%%%%%%%%%%%%%%%%%%%%%
\section{Discrete Spectrum}

\begin{theorem}
If $\lambda>1/2$, the discrete spectrum consists of
\[
E_n = D_e - \frac{\hbar^2 a^2}{2m}(\lambda-n-\tfrac12)^2,
\quad n=0,\dots,n_{\max},
\]
with
\[
n_{\max}=\lfloor\lambda-\tfrac12\rfloor.
\]
\end{theorem}

Eigenfunctions:
\[
\psi_n(x)
=
N_n\, y^{\lambda-n-\frac12}
e^{-y/2}
L_n^{(2\lambda-2n-1)}(y),
\quad
y=2\lambda e^{-a(x-x_0)}.
\]

%%%%%%%%%%%%%%%%%%%%%%%%%%%%%%%%%%%%%%%%%%%%%%%%%%%%%%%%%%%%
\section{Continuous Spectrum}

\begin{theorem}
The essential spectrum is
\[
\sigma_{\text{ess}}(H)=[D_e,\infty).
\]
\end{theorem}

Scattering solutions:
\[
\psi_k(x)
=
e^{ikx}
{}_1F_1\!\left(
\lambda+\tfrac12-\frac{ik}{a},
\,1-\frac{2ik}{a};
\,2\lambda e^{-a(x-x_0)}
\right).
\]

Reflection coefficient:
\[
R(k)=
\frac{\Gamma(-2ik/a)\Gamma(\lambda+\tfrac12+ik/a)}
{\Gamma(2ik/a)\Gamma(\lambda+\tfrac12-ik/a)}.
\]

%%%%%%%%%%%%%%%%%%%%%%%%%%%%%%%%%%%%%%%%%%%%%%%%%%%%%%%%%%%%
\section{Resonances}

Poles of $R(k)$ in $\Im k<0$ occur at
\[
\lambda+\tfrac12-\frac{ik}{a}=-n,
\]
giving
\[
k_n=-i a(\lambda+\tfrac12+n).
\]

These define Gamow states satisfying outgoing boundary conditions.

%%%%%%%%%%%%%%%%%%%%%%%%%%%%%%%%%%%%%%%%%%%%%%%%%%%%%%%%%%%%
\section{Hard-Wall Limit $a\to\infty$}

\begin{theorem}
As $a\to\infty$,
\[
H\to
-\frac{\hbar^2}{2m}\frac{d^2}{dx^2}+D_e
\quad\text{on }(x_0,\infty)
\]
with Dirichlet boundary at $x_0$.
\end{theorem}

\begin{proof}[Sketch]
For $x<x_0$, $V\to\infty$.
For $x>x_0$, $V\to D_e$.
Norm-resolvent convergence holds.
\end{proof}

Reflection:
\[
R(k)\to -1.
\]

Discrete spectrum disappears.

%%%%%%%%%%%%%%%%%%%%%%%%%%%%%%%%%%%%%%%%%%%%%%%%%%%%%%%%%%%%
\section{Comparison with P\"oschl--Teller and Scarf}

\subsection{Hyperbolic P\"oschl--Teller}

\[
V(x)=
-\frac{\hbar^2 a^2}{2m}
\frac{\lambda(\lambda-1)}{\cosh^2(ax)}.
\]

Bound states:
\[
E_n=-\frac{\hbar^2 a^2}{2m}(\lambda-1-n)^2.
\]

Large-$a$ limit localizes to delta-type interaction.
Large-$\lambda$ limit produces deep symmetric well.

\subsection{Trigonometric P\"oschl--Teller}

Defined on finite interval.
Hardening produces infinite square well.
Spectrum remains purely discrete.

\subsection{Scarf Potentials}

Scarf I: compact interval ($\mathfrak{su}(2)$).
Scarf II: full line ($\mathfrak{su}(1,1)$).
Large parameter limits deepen wells rather than create boundaries.

\subsection{Geometric Distinction}

\[
\text{Morse: boundary formation}
\qquad
\text{P\"oschl--Teller/Scarf: well deepening}.
\]

%%%%%%%%%%%%%%%%%%%%%%%%%%%%%%%%%%%%%%%%%%%%%%%%%%%%%%%%%%%%
\section{Algebraic Contraction}

Morse and hyperbolic PT realize $\mathfrak{su}(1,1)$.
Trigonometric PT realizes $\mathfrak{su}(2)$.

Hard-wall limit corresponds to contraction:
\[
\mathfrak{su}(1,1)
\longrightarrow
\text{free translation algebra on half-line}.
\]

%%%%%%%%%%%%%%%%%%%%%%%%%%%%%%%%%%%%%%%%%%%%%%%%%%%%%%%%%%%%
\section{Semiclassical Interpretation}

Bohr–Sommerfeld:
\[
\oint p(x)\,dx
=
2\pi\hbar(n+\tfrac12).
\]

As $a\to\infty$:
Turning point collapses to boundary.
Action integral reduces to half-line free motion.

%%%%%%%%%%%%%%%%%%%%%%%%%%%%%%%%%%%%%%%%%%%%%%%%%%%%%%%%%%%%
\appendix

\section{Weyl’s Alternative}

For Sturm–Liouville operator
\[
-\psi''+V\psi=E\psi,
\]
an endpoint is:

\begin{itemize}
\item limit-point if at most one $L^2$ solution,
\item limit-circle if two $L^2$ solutions.
\end{itemize}

Morse is limit-point at both infinities.

%%%%%%%%%%%%%%%%%%%%%%%%%%%%%%%%%%%%%%%%%%%%%%%%%%%%%%%%%%%%
\begin{thebibliography}{99}

\bibitem{ReedSimon}
M. Reed and B. Simon,
\textit{Methods of Modern Mathematical Physics II},
Academic Press (1975).

\bibitem{Teschl}
G. Teschl,
\textit{Mathematical Methods in Quantum Mechanics},
AMS (2009).

\bibitem{Flugge}
S. Fl\"ugge,
\textit{Practical Quantum Mechanics},
Springer (1971).

\bibitem{CooperKhareSukhatme}
F. Cooper, A. Khare, U. Sukhatme,
Supersymmetry in Quantum Mechanics,
Phys. Rep. 251 (1995).

\bibitem{Gendenshtein}
L. Gendenshtein,
JETP Lett. 38 (1983).

\bibitem{Taylor}
J. R. Taylor,
\textit{Scattering Theory},
Wiley (1972).

\end{thebibliography}

\end{document}
